% ===========================================================================
\section{Discussion}
\label{sec:precise-alignment-discussion}
% ===========================================================================

This section interprets the quantitative results presented in \cref{sec:precise-alignment-results}, examining convergence behaviour, the effects of experimental variables (distance, angle), and the comparative strengths and limitations of the LiDAR and stereo modalities. The cross-modal trade-offs are synthesised in \cref{subsec:pa-discussion-cross-modal}.

A methodological caveat applies throughout this discussion. The fitness score and inlier RMSE reported in \cref{sec:precise-alignment-results} are \emph{internal} registration metrics: they measure geometric agreement between the source and target point clouds after alignment, not absolute pose accuracy relative to a ground-truth reference. A high fitness score is necessary but not sufficient for correct alignment; interpretations of RMSE as ``precision'' should be understood in this restricted sense. The absence of external ground truth is discussed further in \cref{subsec:pa-discussion-limitations}.

% ---------------------------------------------------------------------------
\subsection{Convergence Reliability}
\label{subsec:pa-discussion-convergence}
% ---------------------------------------------------------------------------

A pipeline that achieves low RMSE on converging trials but fails silently on a substantial fraction of conditions offers limited operational utility. In this regard, the LiDAR modality shows a clear advantage: all nine conditions converged, with 311 of 315 hypotheses accepted (98.7\% hypothesis-level acceptance). The stereo pipeline converged in seven of nine conditions (78\%).

The unperturbed initial hypothesis (Orig group) converged in all 9 conditions, with a per-hypothesis acceptance rate of 100\%, consistent with the perturbed groups (97--100\%; \cref{tab:pa-results-groups}). Furthermore, when perturbation hypotheses do converge, they reach effectively the same final pose as the initial hypothesis (\cref{tab:pa-results-lidar-sigma}). These results indicate that the perturbation strategy did not rescue any condition that would otherwise have failed; the LiDAR pipeline's convergence advantage over stereo is attributable primarily to the properties of the LiDAR-derived cost surface rather than to the multi-hypothesis architecture.

Two properties of active LiDAR sensing contribute to this well-conditioned cost surface. First, as an active sensor emitting at 905\,nm, the Livox Mid-360 LiDAR generates returns independently of ambient illumination and surface texture, producing consistent point density across conditions. Second, the near-zero within-trial pose standard deviation observed across most conditions (\cref{tab:pa-results-lidar-sigma}) suggests that the GICP cost surface contains a dominant basin sufficiently wide that all 35 perturbation hypotheses fall within it whenever the preprocessing stage produces a well-formed source cloud.

The two stereo-only failures reveal different vulnerability modes. D40A$+$15 produced 87\,994 processed points, a count comparable to the successfully converging D40A$-$15 (87\,772 points) and D40A0 (83\,163 points), yet achieved $F = 0.0$. Importantly, LiDAR's unperturbed initial hypothesis also converged on this condition, demonstrating that the failure is specific to the stereo-derived cost surface rather than an initialisation problem that perturbations could resolve. The stereo pipeline's dependence on photometric contrast for disparity estimation, combined with the geometric configuration at close range and oblique angle, likely produces a cost surface with a narrower or shallower convergence basin than the LiDAR equivalent. D80A$-$15 represents a different failure: only 23\,980 points survived preprocessing (compared to 42\,307 for D80A0 and 29\,981 for D80A$+$15), suggesting that the trunk segmentation stage was overly aggressive at this particular combination of distance and viewing angle, removing geometrically relevant structure and leaving insufficient correspondences for registration.

Although the perturbation strategy did not determine the convergence outcome in any condition tested, it provides a probabilistic safety margin against future conditions where the initial pose may fall outside the convergence basin, which is more likely in field deployment where SLAM drift is larger and vine geometry is more variable than in this controlled laboratory setting. The multi-hypothesis architecture thus functions as insurance rather than as the mechanism responsible for the observed reliability. For the stereo pipeline, the convergence gap is primarily a sensing and cost-surface conditioning problem: introducing perturbations would address the narrow-basin vulnerability demonstrated by D40A$+$15, but would not resolve failures attributable to insufficient point density (D80A$-$15). Improving stereo preprocessing robustness, particularly the trunk segmentation and depth filtering stages, would yield more direct reliability gains than adding perturbation hypotheses.

% ---------------------------------------------------------------------------
\subsection{Distance Effects on Registration Quality}
\label{subsec:pa-discussion-distance}
% ---------------------------------------------------------------------------

A counter-intuitive finding emerges from the distance analysis: the D80 group (80\,cm) achieves the lowest median inlier RMSE for both modalities (23.77\,mm for LiDAR and 12.45\,mm for stereo), outperforming both the closer D40 group (34.66\,mm LiDAR, 16.77\,mm stereo) and the D60 group (34.29\,mm LiDAR, 16.67\,mm stereo). An intuitive expectation is that closer sensing yields higher precision due to higher point density on the target structure; however, the data consistently indicate that moderate distance produces superior registration quality.

For the LiDAR modality, the geometric explanation centres on the spatial extent of structure captured within the sensor's field of view and the resulting constraint on rotational degrees of freedom. At 40\,cm, the Livox Mid-360's field of view captures only a compact segment of the vine structure: a short length of trunk and its immediate lateral canes. Although this close-range cloud is not geometrically featureless, its limited spatial extent provides a short effective lever arm for constraining rotation. A small angular misalignment applied to a spatially compact cloud produces only small positional residuals at the cloud's edges, leaving the ICP cost surface relatively flat in rotational directions and permitting the solver to terminate with residual angular error that elevates RMSE. At 80\,cm, the wider field of view encompasses a substantially larger scene: multiple vine structures, trellis posts, and guide wires in addition to the target vine's branching geometry. This extended spatial baseline means that the same rotational error produces proportionally larger residuals at distant points, sharpening the cost surface and constraining the solver to a tighter minimum. The effect is analogous to increasing baseline length in triangulation: longer baselines yield better angular resolution. The IQR data support this interpretation: D80 conditions exhibit near-zero IQR in both fitness and RMSE (\cref{tab:pa-results-lidar-per-condition}), indicating that the solver converges to a single well-defined minimum, whereas D40 conditions show slightly elevated IQR values consistent with residual rotational ambiguity from the shorter spatial baseline.

The stereo modality exhibits the same distance-RMSE trend (D80 $<$ D60 $<$ D40), though the absolute RMSE values are approximately half those of LiDAR across all distance groups. The stereo pipeline applies a near-field depth clipping threshold ($z_\text{far} = 1.0$\,m), retaining only high-confidence depth estimates within the stereo baseline's optimal operating range. At D80, this clipping window captures the vine structure at a distance where the stereo disparity estimation is geometrically well-conditioned: the baseline-to-depth ratio yields sub-pixel disparity that translates to millimetre-level depth precision, consistent with the quadratic depth-error relationship for triangulation-based systems~\cite{hartleyMultipleViewGeometry2004}. At D40, the same clipping window may include points at the extreme edges of the depth range where stereo triangulation uncertainty increases, and the very close range may introduce depth-edge artefacts from the narrow stereo baseline. The net effect is a modest but consistent RMSE advantage for the D80 group.

The practical recommendation emerging from these results is an optimal operating range of 60--80\,cm for both modalities. Within this range, convergence is reliable ($\geq 97\%$ for LiDAR, 67--100\% for stereo) and RMSE is minimised. The results also indicate that registration degrades gracefully as distance decreases below 80\,cm (higher RMSE but maintained convergence), suggesting that approaching closer than the optimal distance produces suboptimal but usable alignments. For context, prior work on LiDAR-based vine reconstruction by \citet{morenoOnGroundVineyardReconstruction2020} used a mobile mapping platform with the sensor positioned at approximately 1\,m from the crop row, though direct RMSE comparison is complicated by differences in registration targets (full canopy vs.\ isolated trunk). The sub-25\,mm LiDAR RMSE and sub-13\,mm stereo RMSE achieved at D80 place registration precision within the same order of magnitude as the target structures (5--15\,mm cane diameter), satisfying the centimetre-level accuracy goal stated in \cref{sec:research-objectives}.

% ---------------------------------------------------------------------------
\subsection{Approach Angle Effects}
\label{subsec:pa-discussion-angle}
% ---------------------------------------------------------------------------

The experimental design tested three approach angles ($-15^{\circ}$, $0^{\circ}$, $+15^{\circ}$) to assess whether oblique sensor-to-vine geometry degrades registration performance. For the LiDAR modality, the answer is negative within the tested range: both the A$-$15 and A$+$15 groups achieve 100\% hypothesis acceptance, with RMSE values (34.66\,mm and 34.28\,mm respectively) indistinguishable from the A0 group (33.25\,mm). The A0 group achieves 96\% acceptance (101/105 hypotheses), with the slight reduction attributable to D80A0 (32/35). The LiDAR pipeline is thus effectively invariant to $\pm 15^{\circ}$ approach angle variation, which is an important property for operational robustness given that a mobile platform navigating between vine rows cannot guarantee a perfectly perpendicular approach.

The stereo modality shows convergence rates of 67\% for A$-$15 and A$+$15, and 100\% for A0. D40A$+$15 failed while D40A$-$15 succeeded under nominally symmetric angular conditions. This asymmetry suggests an interaction between approach angle and the specific vine geometry at close range: a $+15^{\circ}$ approach from 40\,cm may position the cameras such that the trunk is partially occluded by a lateral cane or the initial pose estimate is offset into a region where ICP diverges.

The $\pm 15^{\circ}$ range tested here was chosen to represent a conservative estimate of lateral (yaw) approach angle variation for a mobile manipulator operating in vineyard rows. In practice, imprecise stopping (slightly early or late along the row) introduces a yaw offset between the sensor boresight and the vine centre. The stereo pipeline's relatively low convergence on angled approaches limits its viability as a standalone alignment modality; whereas the LiDAR pipeline's invariance demonstrates a more operationally robust modality.

Two additional limitations constrain the generality of the angular analysis. First, only lateral (yaw) deviations were tested. Vineyard terrain is frequently uneven, causing the robot chassis to pitch or roll relative to the vine. Vertical angular deviations could alter the point cloud overlap with the reference model in ways not captured by the lateral-only design. Second, even where the alignment pipeline tolerates $\pm 15^{\circ}$ approach angles, this does not imply that the downstream pruning operation is viable at such offsets. The robotic manipulator's reachability, joint limits, and tool-vine geometry impose additional constraints on acceptable approach angles that were not evaluated in this work. Characterising the manipulator's angular tolerance and its interaction with alignment accuracy would require integrated system-level testing beyond the scope of this evaluation.

% ---------------------------------------------------------------------------
\subsection{Perturbation Strategy Effectiveness}
\label{subsec:pa-discussion-perturbation}
% ---------------------------------------------------------------------------

The multi-hypothesis perturbation strategy generates 35 initial pose hypotheses per trial, distributed across five groups (Orig, A--E) that sample translational and rotational offsets around the SLAM-provided initial pose. Two lines of evidence indicate that these hypotheses all fall within a single dominant convergence basin. First, per-group acceptance rates are uniformly high---97\% (Group~C) to 100\% (Orig, A, B)---with median fitness and RMSE values statistically indistinguishable across groups (\cref{tab:pa-results-groups}). If local minima existed within the perturbation neighbourhood, some hypotheses would become trapped, producing disparity in convergence rates; this is not observed. Second, the within-trial convergence analysis (\cref{tab:pa-results-lidar-sigma}) shows that accepted hypotheses converge to effectively the same final pose: D80A0 achieves near-zero standard deviation across all transformation components ($\sigma < 0.1$\,mm translation, $< 0.01^{\circ}$ rotation), and D40A0, D60A$-$15, and D60A0 show similarly tight clustering. The Orig (unperturbed) group converged in every condition where any hypothesis converged (9/9), confirming that the SLAM-provided initial pose already falls within the basin.

The perturbation strategy therefore provides \emph{redundancy} rather than \emph{diversity}: its value in these experiments was not in rescuing convergence but in providing a probabilistic safety margin for future conditions where the nominal pose may fall outside the basin. With per-hypothesis convergence probability $p$, 35 independent hypotheses achieve trial-level convergence probability $1 - (1-p)^{35}$, which exceeds 99\% for any $p > 13\%$. In the present evaluation, this margin was never exercised.

The exceptions to single-basin convergence are informative. D80A$-$15 and D80A$+$15 exhibit elevated roll standard deviation ($\sigma_\phi > 6^{\circ}$), and D40A$-$15 and D40A$+$15 show moderate roll variability ($\sigma_\phi \approx 4^{\circ}$). These conditions are exclusively the angled approaches, and two factors likely contribute to the elevated roll variability. First, oblique viewing angles reduce the geometric constraint available to the ICP solver in the roll direction: when the sensor views the vine off-axis, fewer surface normals are oriented to penalise roll misalignment, producing a flatter cost surface in that degree of freedom. Second, the manually-measured initial poses for angled conditions carry higher angular uncertainty than the straight-on approaches, because positioning the robot at a precise $15^{\circ}$ offset is inherently more difficult to measure than a perpendicular alignment. This larger initial angular spread, combined with the weaker roll constraint, causes hypotheses to settle at slightly different roll values rather than converging to a single sharp minimum. The roll variability does not affect the translational components ($\sigma_{t_x}$, $\sigma_{t_y}$, $\sigma_{t_z}$ remain low for these conditions), confirming that the positional alignment is well-determined even when the rotational alignment around the trunk axis is less precisely constrained.

For operational deployment, the redundancy interpretation suggests that the number of hypotheses could be substantially reduced without sacrificing reliability. If 35 hypotheses converge to the same solution in 8 of 9 conditions, then 7--10 hypotheses would provide equivalent coverage with a threefold to fivefold reduction in computation time. The primary exception is D80A0, where only 32 of 35 hypotheses were accepted; it still shows a 91\% per-hypothesis acceptance rate, for which 10 hypotheses would yield $>$99.99\% trial-level convergence probability. Such a reduction would bring the LiDAR pipeline's timing from $\sim$46\,s down to $\sim$12--15\,s, approaching the stereo pipeline's latency and substantially improving real-time feasibility without requiring algorithmic parallelisation.

% ---------------------------------------------------------------------------
\subsection{Cross-Modal Trade-offs}
\label{subsec:pa-discussion-cross-modal}
% ---------------------------------------------------------------------------

The matched-condition comparison (\cref{tab:pa-results-cross-modal}) summarises the central finding of this evaluation: the LiDAR and stereo modalities occupy complementary positions in the reliability--precision--speed trade-off space. Neither modality dominates the other across all axes of performance, and the choice between them (or the strategy for combining them) depends on the priorities of the downstream operations.

On the seven conditions where both pipelines converged, stereo achieves a mean inlier RMSE of 15.51\,mm compared to 31.40\,mm for LiDAR, a factor-of-two precision advantage. This difference is not primarily attributable to the ICP algorithm (both pipelines use the same GICP solver with identical convergence parameters) but to the quality of the input point clouds. At close range, stereo imaging produces denser and more spatially uniform point clouds than mechanical LiDAR, which remains comparatively sparse even after prolonged scan accumulation~\cite{stenEnhancingOffRoadTopography2025}. The stereo pipeline yields a compact, high-fidelity cloud of 68\,k--92\,k points from a single frame pair, with near-uniform surface sampling concentrated on the vine structure. The LiDAR pipeline accumulates 97\,k--133\,k points over 25\,s, yet these are distributed across a wider spatial extent with less uniform surface coverage. The higher point density and surface completeness of the stereo cloud provides GICP with a more tightly constrained correspondence set, contributing to the observed increased precision. RMSE as reported here measures the \emph{internal consistency} of inlier correspondences; the mean distance between matched source-target point pairs after alignment, rather than absolute positional accuracy relative to a ground-truth pose. A low RMSE indicates tight geometric agreement between the registered clouds, which is a close proxy condition for pose accuracy but not a comprehensive one.

The reliability advantage of LiDAR is clear: 100\% condition-level convergence versus 78\% for stereo. LiDAR converged in two conditions where stereo failed (D40A$+$15 and D80A$-$15), while no condition that failed for LiDAR succeeded for stereo. As discussed in \cref{subsec:pa-discussion-convergence}, this asymmetry is primarily attributable to the properties of active sensing (illumination independence, consistent point density, and the resulting well-conditioned cost surface) rather than to the multi-hypothesis architecture, which provided redundancy but did not rescue any condition that the unperturbed initial hypothesis would have failed. Stereo, as a passive modality, depends on photometric contrast for disparity estimation; in a laboratory environment this is not limiting, but in field conditions where lighting varies and surfaces may be wet or uniformly coloured, the reliability gap could widen further. The result that LiDAR succeeds everywhere stereo succeeds (but not vice versa) suggests that stereo failures are sensing or preprocessing problems rather than fundamental geometric insufficiency.

The timing differential is substantial. The stereo pipeline completes in $14.18 \pm 2.59$\,s with no preceding accumulation interval, whereas the LiDAR pipeline requires $46.0 \pm 3.8$\,s of processing plus a 25\,s scan accumulation phase (though accumulation can overlap with prior-vine pruning). In pruning systems where the total per-vine cycle is on the order of 100--120\,s~\cite{botterillRobotSystemPruning2017, hizatateDevelopmentRobotPruning2025}, the stereo pipeline's 14\,s alignment latency represents a minor fraction of cycle time, whereas the LiDAR pipeline's 71\,s combined accumulation and processing time becomes the dominant contributor to cycle duration. However, as discussed in \cref{subsec:pa-discussion-perturbation}, reducing the hypothesis count from 35 to 7--10 would bring LiDAR processing from $\sim$46\,s down to $\sim$12--15\,s, yielding a combined latency of $\sim$37--40\,s (including the 25\,s scan accumulation) that is more readily accommodated within the pruning cycle.

The cross-modal comparison also validates the two-stage localisation architecture proposed in this thesis. The ICP corrections reported in \cref{tab:pa-results-lidar-icp-correction,tab:pa-results-stereo-transforms} range from 46.6\,mm to 165.3\,mm in translational displacement and from $0.13^{\circ}$ to $11.5^{\circ}$ in rotational correction. These magnitudes are consistent with the SLAM system's (GLIM-GNSS) absolute trajectory error of 71--174\,mm reported in \cref{chap:SLAM}, confirming that the coarse localisation delivered by SLAM places the robot within ICP's convergence basin. The perturbation strategy samples $\pm 50$\,mm translation and $\pm 15^{\circ}$ rotation around the SLAM pose, and achieves 98.7\% hypothesis-level convergence, which is evidence that the SLAM accuracy is sufficient for the precise-alignment stage to function reliably. The two-stage architecture thus combines capabilities that neither stage provides independently: SLAM provides global consistency and centimetre-level initial pose, while ICP refines this to millimetre-level alignment against the reference model.

% ---------------------------------------------------------------------------
\subsection{Pipeline Timing and Real-Time Feasibility}
\label{subsec:pa-discussion-timing}
% ---------------------------------------------------------------------------

The concept of ``real-time'' in the context of robotic vine pruning is defined not by absolute latency thresholds but by alignment latency as a proportion of the total per-vine cycle. Published pruning systems report per-vine cycle times of approximately 120\,s~\cite{botterillRobotSystemPruning2017} and 107\,s~\cite{hizatateDevelopmentRobotPruning2025}, encompassing perception, planning, and manipulator execution. Alignment is necessarily sequential---motion planning cannot begin until the aligned pose is available---so its latency adds directly to cycle time. Under this framing, an alignment latency of 10--15\,s ($\sim$10--14\% of the per-vine cycle) has minimal impact on throughput, whereas latency exceeding 60\,s would increase cycle time by more than 50\% and become the dominant bottleneck.

The stereo pipeline, with a mean total latency of $14.18 \pm 2.59$\,s, already satisfies this real-time criterion. Its timing decomposition (\cref{tab:pa-results-stereo-timing}) reveals that the dominant components are ICP alignment ($7.41 \pm 2.56$\,s) and dual-camera combination ($4.60 \pm 0.16$\,s), with GPU-accelerated depth estimation via FoundationStereo contributing only 1.2\,s. The combination stage ($4.60 \pm 0.16$\,s) is largely a synchronisation cost: each camera's depth estimation and filtering must complete sequentially before the trivial point cloud merge can proceed. Since the GPU-accelerated stereo inference serialises access to the GPU, the two cameras cannot be processed fully in parallel without additional hardware resources. The ICP latency is determined by cloud size and convergence speed, both of which are already favourable for the stereo pipeline's relatively compact clouds (68\,000--92\,000 processed points). No architectural changes are required for the stereo pipeline to operate within the pruning cycle's time budget.

The LiDAR pipeline's current latency of $46.0 \pm 3.8$\,s (\cref{tab:pa-results-lidar-timing}), plus 25\,s scan accumulation, substantially exceeds the real-time threshold. However, the timing structure illustrates that the majority of the latency is from the sequential execution of ICP, rather than a fundamental computational limit. These hypotheses are \emph{embarrassingly parallel}: they share the same source and target clouds, differ only in their initial transformation, and produce independent results. Since the hypotheses require no inter-process communication, concurrent execution on parallel hardware would reduce the registration phase from $\sim$34\,s to the cost of a single hypothesis ($\sim$1\,s). Combined with the 0.71\,s preprocessing overhead, the parallelised ICP LiDAR pipeline would achieve $\sim$2\,s active computation time, which would be faster than the stereo pipeline. The 25\,s accumulation interval can be overlapped with the preceding vine's pruning execution, effectively removing it from the critical path. With parallelisation and accumulation overlap, the LiDAR pipeline's effective per-vine latency would approach 2--3\,s, well within real-time constraints. the hypothesis reduction discussed in \cref{subsec:pa-discussion-perturbation} (from 35 to 7--10 hypotheses) would achieve a similar latency improvement ($\sim$10--12\,s) without requiring parallel hardware, representing a path to real-time operation that requires minimal engineering effort.

% ---------------------------------------------------------------------------
\subsection{Limitations}
\label{subsec:pa-discussion-limitations}
% ---------------------------------------------------------------------------

The most significant limitation of this evaluation is the absence of an external six-degree-of-freedom ground-truth measurement. All reported metrics (fitness score and inlier RMSE) are \emph{internal consistency} measures: they quantify how well the source and target point clouds agree geometrically after registration, but cannot determine whether the recovered pose is correct in an absolute sense. A registration that achieves high fitness and low RMSE could, in principle, converge to a self-consistent but incorrect local minimum, particularly if both clouds share systematic biases (e.g., both derived from the same sensor platform with an uncalibrated extrinsic offset). The visual alignment overlays (\cref{fig:pa-lidar-alignment-examples,fig:pa-stereo-converged-overlay}) provide qualitative evidence that the registrations are geometrically plausible, but qualitative assessment cannot substitute for quantitative validation against an independent pose reference such as a motion capture system or precision-surveyed fiducial markers. Future work should incorporate such ground truth to establish absolute pose accuracy bounds and determine whether the RMSE values reported here correspond to comparable levels of absolute error.

The evaluation was conducted on two dormant vine specimens in a controlled laboratory environment. While this approach eliminates confounding factors and permits precise control of experimental variables, it leaves open the question of generalisation. Field conditions introduce multiple sources of variability absent from the laboratory: wind-induced motion of canes and foliage (particularly relevant for non-dormant canopy), varying vine structure across cultivars, training systems, and ages, canopy interference from adjacent rows, variable ambient lighting affecting stereo depth estimation, and surface contamination (mud, water droplets) affecting LiDAR reflectance. Each of these factors could independently degrade registration performance, and their combined effect in field conditions is unknown. 

Conversely, the laboratory environment under-represents the geometric information available to the LiDAR pipeline in a vineyard setting. In these experiments, the accumulated LiDAR cloud was cropped to a narrow angular aperture containing only the immediate test vine, because the surrounding laboratory environment is not representative of the operational environment. In a vineyard row, the full 360° LiDAR scan would capture multiple adjacent vines simultaneously, each of which could be registered against its corresponding reference model. Aligning against several spatially distributed reference structures would provide substantially more geometric constraint than a single-vine registration, likely improving both convergence reliability and pose precision for the LiDAR modality. The single-vine, narrow-aperture evaluation presented here therefore represents a conservative assessment of LiDAR alignment performance relative to field deployment.

The experimental design employs a $3 \times 3$ factorial with a single observation per condition ($n = 1$). This design is sufficient to characterise broad performance trends (distance effects, modality differences) but insufficient for rigorous statistical inference. Without replicates, it is impossible to distinguish genuine condition effects from trial-specific noise. For example, the D40A$+$15 stereo failure could reflect a systematic vulnerability at that condition or a stochastic event that would not recur on a second trial. The absence of repeated measurements also precludes assessment of measurement repeatability: that is, the extent to which the same condition produces the same alignment result across multiple independent trials. For a pruning system that must reliably position a cutting tool within millimetres of a target, repeatability is as important as accuracy for operational deployment, and this evaluation provides no information on that axis.

The condition-specific stereo failures (D40A$+$15 and D80A$-$15) may be partially specific to the test vine's geometry and the chosen preprocessing parameters. The D40A$+$15 stereo failure may reflect a geometric interaction between the particular arrangement of canes on this vine and the camera viewpoint at that condition, rather than a general vulnerability of the stereo pipeline at 40\,cm with $+15^{\circ}$ approach angle. Similarly, the D80A$-$15 failure appears linked to overly aggressive trunk segmentation at that viewing angle. Characterising the robustness of these failure boundaries would require testing across multiple vine specimens in several different vineyards, which is left as future work.

Finally, the preprocessing pipeline represents a confounding factor in attributing performance characteristics to the sensing modality versus the algorithmic choices. Failures attributed to ``the LiDAR pipeline'' conflate the Livox Mid-360's sensing capabilities with the specific Patchwork++ parameterisation, the statistical outlier rejection thresholds, and the voxel downsampling resolution; similarly, ``stereo pipeline'' failures conflate FoundationStereo's depth estimation with the DBSCAN clustering parameters, morphological closing kernel size, and the $z_\text{far}$ clipping threshold. A different parameterisation of these preprocessing stages might shift the convergence boundaries or alter the RMSE characteristics. Disentangling the contributions of sensing physics from algorithmic preprocessing choices would require ablation studies that systematically vary preprocessing parameters while holding the sensing modality constant, an analysis beyond the scope of this evaluation but necessary for informed pipeline optimisation in future deployments.
